\documentclass[openany,fontset=fandol]{ctexbook}
\usepackage{amsmath,amssymb}
\usepackage[unicode,hidelinks]{hyperref}
\setlength{\emergencystretch}{2em}
\begin{document}
\title{代数几何基础第五卷预备笔记\\普通奇异零点与几何正规性}
\author{据英译编辑稿翻译，参校法文扫描}
\date{2026年11月20日}
\maketitle
\chapter*{范围与接续}
本批译出英译第67页“Remark. With the mentioned conditions”起的连续正文，至第68页仿射平面例结束，对读法文扫描第102--103页。下一起点为英译第68页引入可逆模截面的类比图之句；对应法文第103页手写引图句。本稿是独立增补，未并入第十七版。
\chapter*{备注：普通奇异零点（续）}
\noindent\textbf{备注。}
在关于$X,Y,S,\varphi$的上述条件下，由于
\[
U=V(\varphi)_{\mathrm{sing}}-V(\varphi)_{\mathrm{sup\,sing}}
\]
在$Y$上非分歧，它在$Y$上的纤维是离散的。这意味着，对每个$y\in Y$，$X_y$中$\varphi_y$的每一个普通奇异零点在奇异零点集合中都是孤立的，也就是说，它在$V(\varphi_y)/k(y)$的非光滑点集合中是孤立的。

于是立即得出：若$d\ge 2$，它就是$V_y$在$k(y)$上的一个几何正规点（因为无论如何，$V_y$都是 Cohen--Macaulay 的）。\footnote{原文这一句附有“p.c.x.”，法文扫描为手写插入；此处保留缩写，不擅自展开。}

在定义1的条件下，可以更一般地证明：若$x$是一个普通奇异零点，则$x$是$V(\varphi)$的一个孤立奇点，其含义是，$x$在$V(\varphi)$中的每一个严格泛化点$x'$都是$V(\varphi)$的正则点。因此，若$V(\varphi)$的奇异点集是闭的，例如$X$是优良概形时，$x$也是该奇异点集的一个孤立点。（若求周全，应当把这个结论收入第四章，但第16节似乎并不是合适的位置。放在哪里呢？）

由此可见，若$X$在域$k$上光滑，而$x$是一个几何普通奇异零点，则$x$是$V(\varphi)$的非光滑点集合中的孤立点；因此，若$\dim_x(X)\ge 2$，它便是几何正规的。（但若$\dim_x(X)=1$，则必定不正规；例如可看典型情形$\varphi=xy$，其中$x,y$是仿射平面的坐标函数。）
\chapter*{译校说明（非原文）}
\paragraph{措辞及范围。}
法文第102页在孤立性句中手写补入“即”的意思，与英译一致。第103页所记“每个泛化$x'$”按此处论旨指异于$x$的泛化；若包含$x$本身，就会同时断言该奇点正则。英译中疑为排字错误的“generalization of $x^1$ of $x$”据法文读为“$x$的泛化$x'$”。定义1段是原文提出的进一步结论，不是原文已给出证明。此处保留历史术语“普通奇异零点”；特征$2$下它与双线性判别式的关系仍应参照主稿已有校勘。

\paragraph{正规性的维数疑点。}
英译第67--68页与法文第102--103页均保留正文所写的阈值；尤其末段法文也写$\dim_x(X)\ge2$与$\dim_x(X)=1$。不在译文中替作者改动。这里须分清环境$X$和超曲面$V(\varphi)$的维数：若在几何纤维的局部环上，$V(\varphi)$为 Cohen--Macaulay，并且除闭点外局部均正则，则 Cohen--Macaulay 给出$S_2$；当该局部环维数至少为$2$时，所有高度至多$1$的点都不是闭点，故满足$R_1$，由 Serre 正规性判据得正规。这个论证的维数条件是加在超曲面局部环上，不是直接加在环境上。

原文自己的例子说明不能把末段的环境维数阈值照字面作为无条件结论：在仿射平面的原点，环境维数是$2$，但$xy=0$的局部环有两个分支，非整环，因而不是正规局部环；零概形在此处的维数是$1$。首处$d$究竟取零概形维数还是环境相对维数，须按前文具体定义核对，不因末段疑点而擅自代换。此处的维数辨析是译校推断，不是原文订正版。
\end{document}
